\section{Introduction}\label{sec:intro}

Type-based verification provides a compositional way to reason about \emph{safety} and \emph{reachability} properties through \emph{may}- or \emph{must}-style reasoning, respectively.
These properties can be specified directly by the type language, where the typecheck procedure can then be reinterpreted as a logical deduction to prove that behaviors are consistent with these types. In a may-style type system, a typing judgement $\Gamma \vdash e : \tau$ requires the result value of term $e$ \emph{may} lie in the safety specification specified by type $\tau$ under \emph{all} possible environments admitted by context $\Gamma$, which yields the safety guarantee: a well-typed program never goes wrong under all possible environments. We can interpret this in another way: the type context $\Gamma$ is \emph{demonic}, choosing the worst possible environment adversarially. In order to defend against the demonic context, we should make sure that $e$ is consistent with type $\tau$ under all possible environments that the demonic context can choose. Dually, a must-style type system can enable reachability verification by requiring all behaviors admitted by type $\tau$ \emph{must} be realized by term $e$ as long as there \emph{exists} an environment admitted by the context. Then context $\Gamma$ is interpreted in \emph{angelic} style, which smartly chooses the desired environment according to the reachability goal.

May-style safety verification has been studied for years in type-based verification, where refinement type systems are a prominent example. On the other hand, reachability reasoning via must/angelic-style interpretation has recently been studied in program logic\cite{ReverseHoareLogic, OHP19, ZDS23}, and is used to reason about ``incorrectness'' properties of programs, i.e., the existence of witnesses of buggy program behavior that can avoid overwhelming false positives in bug detection. There are relatively few forms of must-style reasoning in type-based verification. Recent refinement type systems~\cite{RefinementTypeSymbolicExecution,RefinementTypeRefutations} provide must-instantiation for type failure via symbolic execution to search for concrete counterexamples. Some must-style type systems focus on specific reachability properties: coverage type~\cite{CoverageType} reasoning focuses on all desirable test inputs that must be produced by a test input generator, while \citet{Ill-types} focus on a necessary reasoning of ``ill-typed must get stuck''.

However, these systems typically can only choose one of may- or must-style reasoning, which blocks many verification tasks that require considering both safety and reachability properties simultaneously. For example, a type judgement $\Gamma \vdash x\;\Code{div}\;y : \tau_{\Code{even}}$ tries to show there exists some assignment of $x$ and $y$ that may make an even division result. Type $\tau_{\Code{even}}$ is interpreted in may-style while context $\Gamma$ is interpreted in angelic style that smartly chooses a desirable assignment of $x$ and $y$. Moreover, the type of division operator $\Code{div}$ also has a safety guard that the divisor $y$ is non-zero. None of the above systems can unify these notions of may/demonic-style and must/angelic-style reasoning in a foundational type theory, which enables both safety and reachability reasoning in a single type system.

\newcommand{\natminus}{\ensuremath{-_{\Code{nat}}}}

Enabling both may-style and must-style reasoning in type-based verification is challenging in functional languages, where functions themselves can have their own context. This means the context can be mixed with angelic and demonic interpretations in arbitrary ways. Consider the following example:
\begin{align*}
    &x{:}\urt{\Nat}{\vnu > 0} \vdash \lambda y.x - y : y{:}\urt{\Nat}{\vnu > 0} \sarr \urt{\Nat}{\vnu > 0} 
    \\&x{:}\ort{\Nat}{\vnu > 0} \vdash \lambda y.x - y : y{:}\urt{\Nat}{\vnu > 0} \sarr \urt{\Nat}{\vnu > 0}
\end{align*}\noindent
We use $\ort{b}{\phi}$ to express that a primitive type $b$ is refined with a qualifier $\phi(\vnu)$ in may and demonic style, and $\urt{b}{\phi}$ for must and angelic style. In this sense, type $y{:}\urt{\Nat}{\vnu > 0} \sarr \urt{\Nat}{\vnu > 0}$ requires that the function can reach any positive number when parameter $y$ is angelically assigned a positive number.
The two judgements have the same term and types, but different interpretations of their contexts. The first judgement simply asks \emph{if the context can assign $x$ and $y$ angelically, can all positive numbers be reached by $x - y$?} The answer is yes. Assume we want to reach a positive number $n$; then the angelic context can choose a desirable environment $[x \mapsto n + 1, y \mapsto 1]$. More interestingly, the second judgement asks: \emph{if context $\Gamma$ can assign $x$ demonically, can function $\lambda y.x - y$ still reach any positive number with an angelically assigned $y$?} The answer is no, since a demonic context can choose a bad environment $[x \mapsto 0]$, which makes the execution result of term $x - y$ reach only $0$, whatever the angelically assigned $y$ is. This means that a foundational type theory should provide a principled way to compose pieces of demonic and angelic context.

Moreover, the functional language also allows a function to \emph{capture} part of the current context, which makes verification of higher-order functions more challenging. In the previous example, we highlight that there are two input contexts for function $\lambda y.x - y$: a ``captured context'' $x{:}\urt{\Nat}{\vnu > 0}$ contains the variables referred to by the function body, while ``parameter context'' $y{:}\urt{\Nat}{\vnu > 0}$ is the context expected by this function during application. This complicated context dependency makes higher-order functions more difficult to reason about.
\begin{align*}
    x{:}\urt{\Nat}{\vnu > 0}, y{:}\urt{\Nat}{\vnu > 0} &\vdash  \lambda f. f\;x : f{:}\tau_f \sarr \urt{\Nat}{\vnu > 0} \quad\judgfail
\end{align*}\noindent
As shown above, the higher-order function $\lambda f. f\;x$ takes a function $f$ that is assumed to have type $\tau_f$ discussed above and directly applies $f$ to $x$. At first glance, it should guarantee reachability of all positive numbers. However, when $f$ is the function $\lambda y.x - y$, it always returns $0$.
\begin{align*}
    &f_1 \doteq \lambda y.x - y \quad 
    \\&\zlet{x}{e_x}{\boxed{(\lambda f. f\;x)\;f_1}} \;\hookrightarrow^*\; \zlet{x}{e_x}{\boxed{x - x}} \;\hookrightarrow^*\; 0 
\end{align*}\noindent
As shown above, we use $\zlet{x}{e_x}{...}$ to express the environment chosen by the angelic context.
The problem is that the ``captured context'' of function $\lambda y.x - y$ contains variable $x$, which creates a hidden dependency: parameter $y$ is expected not to alias with $x$. Moreover, a similar alias problem can happen with the ``result context'' of function application. Consider the following nested-application example:
\begin{align*}
    x{:}\urt{\Nat}{\vnu > 0}, y{:}\urt{\Nat}{\vnu > 0} \vdash \lambda f. f(f\;y) : f{:}\tau_f \sarr \urt{\Nat}{\vnu > 0} \quad \judgfail
\end{align*}\noindent
Although the first application $f\;y$ works as expected, the second application $f(f\;y)$ uses the result of the previous application, which is problematic. Here is a counter-example:
\begin{align*}
    &f_2 \doteq \lambda y.\zif{y = 1}{x}{2} \quad
    \\& \zlet{x}{e_x}{\zlet{y}{e_y}{\boxed{(\lambda f. f\;(f\;y))\;f_2}}} \;\hookrightarrow^*\;
    \\& \zlet{x}{e_x}{\boxed{f_2\;(f_2\;1)}} \;\hookrightarrow^*\;  \quad \text{\textcolor{gray}{(angelic chooses $y = 1$)}}  
    \\& \zlet{x}{e_x}{\boxed{f\;x}} \;\hookrightarrow^*\; \zlet{x}{e_x}{\boxed{\zif{x = 1}{x}{2}}} \;\hookrightarrow^*\; \boxed{1 \text{ or } 2}
\end{align*}\noindent
The function $\lambda y.\zif{y = 1}{x}{2}$ can have type $\tau_f$, since the parameter context can angelically assign value $1$ to $y$ to make the function return $x$. However, the result of the first application then depends on $x$, which makes the nested application fail to reach any value greater than $2$.

In order to define a fundamental type theory that supports both may/demonic-style and must/angelic-style reasoning, we treat these two styles of interpretation as first-class modalities, and support fine-grained mixing in arbitrary ways. Moreover, to solve the context dependency issue, we treat the \emph{capability} of a context to make demonic or angelic choices as \emph{resources} and introduce a \emph{resource algebra} for principled, compositional decomposition and combination of these capabilities. This approach enables us to handle context disjointness in the same way as in bunched implication (e.g., separation logic) and bunched type systems.
We then introduce \emph{context types} for unifying safety and reachability reasoning in a pure functional language.
We allow type contexts to be connected with dependent conjunction ($\Gamma_1 \ctcomma \Gamma_2$) and disjoint conjunction ($\Gamma_1 \ctstar \Gamma_2$). This bunched modal type system enables us to track complex data-flow dependencies as tree-structured contexts and support framing-based local reasoning.
The context types have the normal function type $\ctsarr$ and magic wand type $\ctwand$ to distinguish whether the ``captured context'' is dependent on or disjoint from the ``parameter context''. For example,
\begin{align*}
    &x{:}\urt{\Nat}{\vnu > 0} \vdash \lambda y.x - y : y{:}\urt{\Nat}{\vnu > 0} \ctwand \urt{\Nat}{\vnu > 0}
\end{align*}\noindent
The magic wand above requires parameter $y$ to be independent of the current context (i.e., $x{:}\urt{\Nat}{\vnu > 0}$), which guarantees that $y$ can be chosen angelically regardless of $x$. The result context will always be dependent on the input context for a deterministic language; however, we also generalize our system to support disjoint result context in a functional language allowing nondeterministic behavior.

\paragraph{Contributions.}
We summarize the main technical contributions as follows.
\begin{itemize}[leftmargin=*]
\item We interpret the capability of a context to make angelic or demonic choices as a resource, and introduce a resource algebra for compositional decomposition and combination of this capability.
\item We define a unified modal \emph{context logic} that supports dependent and disjoint mixing of may/demonic and must/angelic modalities.
\item We design \emph{context types}, a bunched modal type system interpreted in our context logic, to support both safety and reachability reasoning for deterministic functional programs. We provide a generalization of context types to a functional language with randomness, where the result context is allowed to be disjoint from the input context.
\item We formalize our core language and type system in Rocq, and provide corresponding helper tactics in an interactive proof assistant to reason about context types.
\end{itemize}





